\section*{Introducción}
\label{sec:introduccion}

\subsection{Motivación}

Desde 2023 Bolivia atraviesa una escasez de divisas que separó de forma
persistente el tipo de cambio oficial, administrado por el Banco Central de
Bolivia (BCB), del precio que el dólar alcanza en el mercado paralelo
(\emph{dólar blue}). Tras el régimen flexible de 2025--2026, el oficial quedó
unificado mientras el paralelo cotiza libremente vía USDT en plataformas P2P
(Binance, Bybit, Bitget) y fintechs (Airtm, Takenos, El Dorado,
Wallbit). Esa \textbf{brecha} es el
principal termómetro de escasez de divisas y de presión cambiaria del país.

\subsection{Alcance}

Este informe documenta, siguiendo la metodología CRISP-DM
\cite{chapman2000crisp}, la construcción de un pipeline de
datos reproducible que:

\begin{itemize}
  \item recolecta y normaliza cotizaciones de la API pública de
        \href{https://www.dolarbluebolivia.click}{dolarbluebolivia.click}
        (37 endpoints) \cite{dolarblue2026api} y de la Gaceta Oficial de
        Bolivia \cite{gaceta2026};
  \item mantiene series de tiempo alineadas blue/oficial (compra/venta);
  \item produce análisis exploratorio reproducible;
  \item ajusta modelos de \textbf{previsión} del nivel del paralelo y de la
        brecha oficial--paralelo (regresión lineal, OLS con errores HAC y
        ARIMA), evaluados con partición temporal y línea base naïve a un
        horizonte de 1~hora y 1~día (Objetivo 1), y estudia la \textbf{relación}
        entre las publicaciones de la Gaceta y la brecha en horizontes de
        horas a días (Objetivo 2);
\end{itemize}

\subsection{Estructura del documento}

La \hyperref[sec:negocio]{sección~\ref*{sec:negocio}} desarrolla la comprensión
del negocio (problema de predicción, Objetivos 1--2, criterios de éxito,
recursos disponibles). La
\hyperref[sec:datos]{sección~\ref*{sec:datos}} describe las fuentes, el
pipeline de recolección y el diccionario de datos. La
\hyperref[sec:analisis]{sección~\ref*{sec:analisis}} presenta el análisis
exploratorio y la \hyperref[sec:preparacion]{sección~\ref*{sec:preparacion}}
la preparación de datos (fase 3) y la
\hyperref[sec:modelado]{sección~\ref*{sec:modelado}} reporta el modelado
(fase 4). La \hyperref[sec:negocio_eval]{sección~\ref*{sec:negocio_eval}}
traduce las métricas a KPIs de negocio y ROI (fase 5). El despliegue y el
monitoreo (fase 6) quedan como trabajo pendiente.
